\section{Conclusion}
\label{sec:conclusion}

In the end, NEO-Watch does what we wanted it to do. The six required features work, and so do
the three optional ones we chose. A user can log in, see the next approaches on the dashboard,
search for an object, follow it in their favourites and get an alert when an object matches their
criteria. The NASA data is also updated every day on its own.

Not everything went as planned, though. Some ideas from the analysis report were dropped or
simplified because we did not have the time, like the validation of user objects by an
administrator. Others changed once we started coding. Logging in now uses tokens instead of sessions, and the
web service ended up with thirty endpoints organised quite differently from the plan. The database also took us longer than expected to get right, and it
delayed the start of several features. Looking back, we should have fixed the schema earlier and
run the whole application together sooner, instead of discovering the integration problems at
the end.

The application also has limits. It only knows the objects that come close to the Earth during
the three years we import, so it is not a complete catalogue, and some favourites only have one
point on their graph. It depends completely on the NASA API, and when NASA has a problem, so do
we. With more time, we would have liked to store the passwords with a safer algorithm such as
bcrypt, bring back the validation of user objects, and test the Streamlit pages automatically.

On our side, this project taught us a lot. It was the first time most of us worked on the same
code base at four, with branches, pull requests and merge conflicts to solve. We also understood
why tests matter. At first they felt like extra work, but our 152 tests found real bugs that we had
never noticed by using the application. For example, an interrupted NASA import used to start
again from scratch instead of resuming where it had stopped. Overall, we are happy with the result, and we learned as much
from what went wrong as from what worked.
